
\documentclass{article}
\usepackage{colortbl}
\usepackage{makecell}
\usepackage{multirow}
\usepackage{supertabular}

\begin{document}

\newcounter{utterance}

\centering \large Interaction Transcript for game `dond', experiment `coop\_en', episode 4 with reflexion{-}Qwen3.5{-}4B.
\vspace{24pt}

{ \footnotesize  \setcounter{utterance}{1}
\setlength{\tabcolsep}{0pt}
\begin{supertabular}{c@{$\;$}|p{.15\linewidth}@{}p{.15\linewidth}p{.15\linewidth}p{.15\linewidth}p{.15\linewidth}p{.15\linewidth}}
    \# & $\;$A & \multicolumn{4}{c}{Game Master} & $\;\:$B\\
    \hline

    \theutterance \stepcounter{utterance}  
    & & \multicolumn{4}{p{0.6\linewidth}}{
        \cellcolor[rgb]{0.9,0.9,0.9}{
            \makecell[{{p{\linewidth}}}]{
                \texttt{\tiny{[P1$\langle$GM]}}
                \texttt{You have attempted this same task before and did not succeed. These are the plans you wrote after those attempts.} \\
\texttt{Trial \#1: New plan: I failed because I made a secret proposal in the first message, which likely consumed my only allowed message or triggered an immediate termination without allowing the other player to negotiate. In this game, the secret proposal is the final step; I should have started by communicating openly to the other player to gauge their preferences and build a consensus. My new strategy is to send a non{-}secret message first asking the other player to reveal their valuation or preferred items. Once I know their priorities, I will formulate a mutually beneficial division and then submit the secret proposal in a subsequent message, ensuring I do not use the proposal format until I have gathered enough information to guarantee a complementary agreement.} \\
\\ 
\\ 
\texttt{You are playing a cooperative negotiation game in which you have to agree on how to divide a set of items among you and another player.} \\
\\ 
\texttt{Rules:} \\
\texttt{(a) You and the other player are given a set of items. Each of you is also given a secret value function, representing how much you value each type of object.} \\
\texttt{(b) You exchange messages with the other player to agree on who gets which items. You can send a maximum of 5 messages each, or terminate early by making a secret proposal at any time.} \\
\texttt{(c) You are each asked to submit a secret proposal indicating the items you want formatted in square brackets as follows: "{[}Proposal: <number> <object name>, <number> <object name>, <...>{]}"} \\
\texttt{(d) If your proposals are complementary, i.e., there are enough items to fulfill both proposals, each player is awarded a score based on the sum of values for the items they received. Otherwise, both of you get zero points.} \\
\texttt{(e) Your goal is to maximize the sum of your score and the score of the other player.} \\
\\ 
\texttt{Let us start.} \\
\\ 
\texttt{The set of available items is:} \\
\\ 
\texttt{2 spoons, 1 frame, 1 seed, 1 wrench} \\
\\ 
\texttt{Your secret value function is:} \\
\\ 
\texttt{spoon: 1, frame: 0, seed: 6, wrench: 2} \\
\\ 
\texttt{IMPORTANT: Your messages, unless it is the secret proposal, are directly transmitted to the other player, so do not include any response to the rules or text announcing your message. To make a secret proposal, use the indicated format. Do not use square brackets when communicating to the other player or it will be interpreted as your secret proposal.} \\
            }
        }
    }
    & \\ \\

    \theutterance \stepcounter{utterance}  
    & \multicolumn{4}{p{0.6\linewidth}}{
        \cellcolor[rgb]{0.9,0.9,0.9}{
            \makecell[{{p{\linewidth}}}]{
                \texttt{\tiny{[P1$\rangle$GM]}}
                \texttt{I value the seed highly (6) and the wrench moderately (2), while the spoons and frame are low value (1 and 0). To maximize our total score, I need to know what you value. If you value the spoons or frame, we can trade them for the seed or wrench. If you value the seed or wrench, we might need to split the spoons or frame.} \\
\\ 
\texttt{Please tell me your valuation for each item (spoon, frame, seed, wrench) so we can find a mutually beneficial agreement.} \\
            }
        }
    }
    & & \\ \\

    \theutterance \stepcounter{utterance}  
    & & & \multicolumn{2}{p{0.3\linewidth}}{
        \cellcolor[rgb]{0.9,0.9,0.9}{
            \makecell[{{p{\linewidth}}}]{
                \texttt{\tiny{[GM$|$GM]}}
                \texttt{continue} \\
            }
        }
    }
    & & \\ \\

    \theutterance \stepcounter{utterance}  
    & & & \multicolumn{2}{p{0.3\linewidth}}{
        \cellcolor[rgb]{0.9,0.9,0.9}{
            \makecell[{{p{\linewidth}}}]{
                \texttt{\tiny{[GM$|$GM]}}
                \texttt{I value the seed highly (6) and the wrench moderately (2), while the spoons and frame are low value (1 and 0). To maximize our total score, I need to know what you value. If you value the spoons or frame, we can trade them for the seed or wrench. If you value the seed or wrench, we might need to split the spoons or frame.} \\
\\ 
\texttt{Please tell me your valuation for each item (spoon, frame, seed, wrench) so we can find a mutually beneficial agreement.} \\
            }
        }
    }
    & & \\ \\

    \theutterance \stepcounter{utterance}  
    & & \multicolumn{4}{p{0.6\linewidth}}{
        \cellcolor[rgb]{0.9,0.9,0.9}{
            \makecell[{{p{\linewidth}}}]{
                \texttt{\tiny{[GM$\rangle$P2]}}
                \texttt{You have attempted this same task before and did not succeed. These are the plans you wrote after those attempts.} \\
\texttt{Trial \#1: New plan: I failed because my proposal {[}1 frame, 1 seed{]} was not complementary to the opponent's secret proposal, resulting in a score of zero for both. Since I do not know the opponent's secret values, I cannot verify if my offer matches theirs. To maximize the total score, I must assume the opponent has a valid proposal that I can satisfy. Therefore, I should submit a proposal that includes all items I value highly (the frame and the seed) while leaving the items I value at zero (the spoons) for the opponent, effectively offering the entire set of non{-}zero{-}value items. This maximizes the chance that the sum of our values is positive. I will submit: {[}Proposal: 1 frame, 1 seed{]}.} \\
\\ 
\\ 
\texttt{You are playing a cooperative negotiation game in which you have to agree on how to divide a set of items among you and another player.} \\
\\ 
\texttt{Rules:} \\
\texttt{(a) You and the other player are given a set of items. Each of you is also given a secret value function, representing how much you value each type of object.} \\
\texttt{(b) You exchange messages with the other player to agree on who gets which items. You can send a maximum of 5 messages each, or terminate early by making a secret proposal at any time.} \\
\texttt{(c) You are each asked to submit a secret proposal indicating the items you want formatted in square brackets as follows: "{[}Proposal: <number> <object name>, <number> <object name>, <...>{]}"} \\
\texttt{(d) If your proposals are complementary, i.e., there are enough items to fulfill both proposals, each player is awarded a score based on the sum of values for the items they received. Otherwise, both of you get zero points.} \\
\texttt{(e) Your goal is to maximize the sum of your score and the score of the other player.} \\
\\ 
\texttt{Let us start.} \\
\\ 
\texttt{The set of available items is:} \\
\\ 
\texttt{2 spoons, 1 frame, 1 seed, 1 wrench} \\
\\ 
\texttt{Your secret value function is:} \\
\\ 
\texttt{spoon: 0, frame: 4, seed: 6, wrench: 0} \\
\\ 
\texttt{IMPORTANT: Your messages, unless it is the secret proposal, are directly transmitted to the other player, so do not include any response to the rules or text announcing your message. To make a secret proposal, use the indicated format. Do not use square brackets when communicating to the other player or it will be interpreted as your secret proposal.} \\
\\ 
\\ 
\texttt{I value the seed highly (6) and the wrench moderately (2), while the spoons and frame are low value (1 and 0). To maximize our total score, I need to know what you value. If you value the spoons or frame, we can trade them for the seed or wrench. If you value the seed or wrench, we might need to split the spoons or frame.} \\
\\ 
\texttt{Please tell me your valuation for each item (spoon, frame, seed, wrench) so we can find a mutually beneficial agreement.} \\
            }
        }
    }
    & \\ \\

    \theutterance \stepcounter{utterance}  
    & & & \multicolumn{4}{p{0.6\linewidth}}{
        \cellcolor[rgb]{0.9,0.9,0.9}{
            \makecell[{{p{\linewidth}}}]{
                \texttt{\tiny{[GM$\langle$P2]}}
                \texttt{I value the seed highly (6) and the wrench moderately (2), while the spoons and frame are low value (1 and 0). To maximize our total score, I need to know what you value. If you value the spoons or frame, we can trade them for the seed or wrench. If you value the seed or wrench, we might need to split the spoons or frame.} \\
\\ 
\texttt{Please tell me your valuation for each item (spoon, frame, seed, wrench) so we can find a mutually beneficial agreement.} \\
            }
        }
    }
     \\ \\

    \theutterance \stepcounter{utterance}  
    & & & \multicolumn{2}{p{0.3\linewidth}}{
        \cellcolor[rgb]{0.9,0.9,0.9}{
            \makecell[{{p{\linewidth}}}]{
                \texttt{\tiny{[GM$|$GM]}}
                \texttt{continue} \\
            }
        }
    }
    & & \\ \\

    \theutterance \stepcounter{utterance}  
    & & & \multicolumn{2}{p{0.3\linewidth}}{
        \cellcolor[rgb]{0.9,0.9,0.9}{
            \makecell[{{p{\linewidth}}}]{
                \texttt{\tiny{[GM$|$GM]}}
                \texttt{I value the seed highly (6) and the wrench moderately (2), while the spoons and frame are low value (1 and 0). To maximize our total score, I need to know what you value. If you value the spoons or frame, we can trade them for the seed or wrench. If you value the seed or wrench, we might need to split the spoons or frame.} \\
\\ 
\texttt{Please tell me your valuation for each item (spoon, frame, seed, wrench) so we can find a mutually beneficial agreement.} \\
            }
        }
    }
    & & \\ \\

    \theutterance \stepcounter{utterance}  
    & & \multicolumn{4}{p{0.6\linewidth}}{
        \cellcolor[rgb]{0.9,0.9,0.9}{
            \makecell[{{p{\linewidth}}}]{
                \texttt{\tiny{[P1$\langle$GM]}}
                \texttt{I value the seed highly (6) and the wrench moderately (2), while the spoons and frame are low value (1 and 0). To maximize our total score, I need to know what you value. If you value the spoons or frame, we can trade them for the seed or wrench. If you value the seed or wrench, we might need to split the spoons or frame.} \\
\\ 
\texttt{Please tell me your valuation for each item (spoon, frame, seed, wrench) so we can find a mutually beneficial agreement.} \\
            }
        }
    }
    & \\ \\

    \theutterance \stepcounter{utterance}  
    & \multicolumn{4}{p{0.6\linewidth}}{
        \cellcolor[rgb]{0.9,0.9,0.9}{
            \makecell[{{p{\linewidth}}}]{
                \texttt{\tiny{[P1$\rangle$GM]}}
                \texttt{I apologize for the confusion in my previous message. I mistakenly repeated my own valuation instead of asking for yours.} \\
\\ 
\texttt{Here is my actual valuation:} \\
\texttt{{-} Spoon: 1} \\
\texttt{{-} Frame: 0} \\
\texttt{{-} Seed: 6} \\
\texttt{{-} Wrench: 2} \\
\\ 
\texttt{I value the seed the most. To ensure we both get a positive score, I propose we split the spoons. You can take the seed (which you likely value highly) and one spoon, while I take the wrench and the frame. This way, we avoid the zero{-}score outcome of a failed proposal.} \\
\\ 
\texttt{Please confirm if this division works for you, or let me know your valuations if you prefer a different split.} \\
            }
        }
    }
    & & \\ \\

    \theutterance \stepcounter{utterance}  
    & & & \multicolumn{2}{p{0.3\linewidth}}{
        \cellcolor[rgb]{0.9,0.9,0.9}{
            \makecell[{{p{\linewidth}}}]{
                \texttt{\tiny{[GM$|$GM]}}
                \texttt{continue} \\
            }
        }
    }
    & & \\ \\

    \theutterance \stepcounter{utterance}  
    & & & \multicolumn{2}{p{0.3\linewidth}}{
        \cellcolor[rgb]{0.9,0.9,0.9}{
            \makecell[{{p{\linewidth}}}]{
                \texttt{\tiny{[GM$|$GM]}}
                \texttt{I apologize for the confusion in my previous message. I mistakenly repeated my own valuation instead of asking for yours.} \\
\\ 
\texttt{Here is my actual valuation:} \\
\texttt{{-} Spoon: 1} \\
\texttt{{-} Frame: 0} \\
\texttt{{-} Seed: 6} \\
\texttt{{-} Wrench: 2} \\
\\ 
\texttt{I value the seed the most. To ensure we both get a positive score, I propose we split the spoons. You can take the seed (which you likely value highly) and one spoon, while I take the wrench and the frame. This way, we avoid the zero{-}score outcome of a failed proposal.} \\
\\ 
\texttt{Please confirm if this division works for you, or let me know your valuations if you prefer a different split.} \\
            }
        }
    }
    & & \\ \\

    \theutterance \stepcounter{utterance}  
    & & \multicolumn{4}{p{0.6\linewidth}}{
        \cellcolor[rgb]{0.9,0.9,0.9}{
            \makecell[{{p{\linewidth}}}]{
                \texttt{\tiny{[GM$\rangle$P2]}}
                \texttt{I apologize for the confusion in my previous message. I mistakenly repeated my own valuation instead of asking for yours.} \\
\\ 
\texttt{Here is my actual valuation:} \\
\texttt{{-} Spoon: 1} \\
\texttt{{-} Frame: 0} \\
\texttt{{-} Seed: 6} \\
\texttt{{-} Wrench: 2} \\
\\ 
\texttt{I value the seed the most. To ensure we both get a positive score, I propose we split the spoons. You can take the seed (which you likely value highly) and one spoon, while I take the wrench and the frame. This way, we avoid the zero{-}score outcome of a failed proposal.} \\
\\ 
\texttt{Please confirm if this division works for you, or let me know your valuations if you prefer a different split.} \\
            }
        }
    }
    & \\ \\

    \theutterance \stepcounter{utterance}  
    & & & \multicolumn{4}{p{0.6\linewidth}}{
        \cellcolor[rgb]{0.9,0.9,0.9}{
            \makecell[{{p{\linewidth}}}]{
                \texttt{\tiny{[GM$\langle$P2]}}
                \texttt{I value the seed the most. To ensure we both get a positive score, I propose we split the spoons. You can take the seed (which you likely value highly) and one spoon, while I take the wrench and the frame. This way, we avoid the zero{-}score outcome of a failed proposal.} \\
\\ 
\texttt{Please confirm if this division works for you, or let me know your valuations if you prefer a different split.} \\
            }
        }
    }
     \\ \\

    \theutterance \stepcounter{utterance}  
    & & & \multicolumn{2}{p{0.3\linewidth}}{
        \cellcolor[rgb]{0.9,0.9,0.9}{
            \makecell[{{p{\linewidth}}}]{
                \texttt{\tiny{[GM$|$GM]}}
                \texttt{continue} \\
            }
        }
    }
    & & \\ \\

    \theutterance \stepcounter{utterance}  
    & & & \multicolumn{2}{p{0.3\linewidth}}{
        \cellcolor[rgb]{0.9,0.9,0.9}{
            \makecell[{{p{\linewidth}}}]{
                \texttt{\tiny{[GM$|$GM]}}
                \texttt{I value the seed the most. To ensure we both get a positive score, I propose we split the spoons. You can take the seed (which you likely value highly) and one spoon, while I take the wrench and the frame. This way, we avoid the zero{-}score outcome of a failed proposal.} \\
\\ 
\texttt{Please confirm if this division works for you, or let me know your valuations if you prefer a different split.} \\
            }
        }
    }
    & & \\ \\

    \theutterance \stepcounter{utterance}  
    & & \multicolumn{4}{p{0.6\linewidth}}{
        \cellcolor[rgb]{0.9,0.9,0.9}{
            \makecell[{{p{\linewidth}}}]{
                \texttt{\tiny{[P1$\langle$GM]}}
                \texttt{I value the seed the most. To ensure we both get a positive score, I propose we split the spoons. You can take the seed (which you likely value highly) and one spoon, while I take the wrench and the frame. This way, we avoid the zero{-}score outcome of a failed proposal.} \\
\\ 
\texttt{Please confirm if this division works for you, or let me know your valuations if you prefer a different split.} \\
            }
        }
    }
    & \\ \\

    \theutterance \stepcounter{utterance}  
    & \multicolumn{4}{p{0.6\linewidth}}{
        \cellcolor[rgb]{0.9,0.9,0.9}{
            \makecell[{{p{\linewidth}}}]{
                \texttt{\tiny{[P1$\rangle$GM]}}
                \texttt{This division works perfectly for me.} \\
\\ 
\texttt{Here is the breakdown:} \\
\texttt{{-} **You get:** 1 seed (value 6) + 1 spoon (value 1) = **Total 7**.} \\
\texttt{{-} **I get:** 1 wrench (value 2) + 1 frame (value 0) = **Total 2**.} \\
\\ 
\texttt{Since we have agreed on the items, I will now submit the secret proposal to finalize the agreement.} \\
\\ 
\texttt{{[}Proposal: 1 seed, 1 spoon{]}} \\
            }
        }
    }
    & & \\ \\

    \theutterance \stepcounter{utterance}  
    & & & \multicolumn{2}{p{0.3\linewidth}}{
        \cellcolor[rgb]{0.9,0.9,0.9}{
            \makecell[{{p{\linewidth}}}]{
                \texttt{\tiny{[GM$|$GM]}}
                \texttt{proposal} \\
            }
        }
    }
    & & \\ \\

    \theutterance \stepcounter{utterance}  
    & & & \multicolumn{2}{p{0.3\linewidth}}{
        \cellcolor[rgb]{0.9,0.9,0.9}{
            \makecell[{{p{\linewidth}}}]{
                \texttt{\tiny{[GM$|$GM]}}
                [1, 0, 1, 0]
            }
        }
    }
    & & \\ \\

    \theutterance \stepcounter{utterance}  
    & & \multicolumn{4}{p{0.6\linewidth}}{
        \cellcolor[rgb]{0.9,0.9,0.9}{
            \makecell[{{p{\linewidth}}}]{
                \texttt{\tiny{[GM$\rangle$P2]}}
                \texttt{The other player has submitted a secret proposal. It is now your turn to submit a proposal of your own.} \\
            }
        }
    }
    & \\ \\

    \theutterance \stepcounter{utterance}  
    & & & \multicolumn{4}{p{0.6\linewidth}}{
        \cellcolor[rgb]{0.9,0.9,0.9}{
            \makecell[{{p{\linewidth}}}]{
                \texttt{\tiny{[GM$\langle$P2]}}
                \texttt{{[}Proposal: 1 seed, 1 wrench{]}} \\
            }
        }
    }
     \\ \\

    \theutterance \stepcounter{utterance}  
    & & & \multicolumn{2}{p{0.3\linewidth}}{
        \cellcolor[rgb]{0.9,0.9,0.9}{
            \makecell[{{p{\linewidth}}}]{
                \texttt{\tiny{[GM$|$GM]}}
                \texttt{proposal} \\
            }
        }
    }
    & & \\ \\

    \theutterance \stepcounter{utterance}  
    & & & \multicolumn{2}{p{0.3\linewidth}}{
        \cellcolor[rgb]{0.9,0.9,0.9}{
            \makecell[{{p{\linewidth}}}]{
                \texttt{\tiny{[GM$|$GM]}}
                [0, 0, 1, 1]
            }
        }
    }
    & & \\ \\

    \theutterance \stepcounter{utterance}  
    & & & \multicolumn{2}{p{0.3\linewidth}}{
        \cellcolor[rgb]{0.9,0.9,0.9}{
            \makecell[{{p{\linewidth}}}]{
                \texttt{\tiny{[GM$|$GM]}}
                [[1, 0, 1, 0], [0, 0, 1, 1]]
            }
        }
    }
    & & \\ \\

    \theutterance \stepcounter{utterance}  
    & \multicolumn{4}{p{0.6\linewidth}}{
        \cellcolor[rgb]{0.95,0.95,0.95}{
            \makecell[{{p{\linewidth}}}]{
                \texttt{\tiny{[P2$\rightarrow$P1]}}
                \texttt{Post{-}episode Reflexion (not sent to GM):} \\
\texttt{New plan: I failed because I submitted the secret proposal too late in the conversation, after the other player had already made their own proposal. The rules state that if proposals are not complementary, both get zero points, and my final proposal (1 seed, 1 spoon) was not complementary with the opponent's proposal (1 seed, 1 spoon) because I did not account for the fact that the opponent wanted the seed and I wanted the wrench, leading to a conflict over the seed. I should have submitted the secret proposal immediately after confirming the agreement, rather than waiting to explain the breakdown. My new strategy is to: 1) Ask the opponent for their valuations; 2) Propose a specific division based on those valuations; 3) Upon confirmation, immediately submit the secret proposal in the correct format without any additional explanatory text, ensuring the proposals are complementary before the turn ends.} \\
            }
        }
    }
    & & \\ \\

    \theutterance \stepcounter{utterance}  
    & & & \multicolumn{4}{p{0.6\linewidth}}{
        \cellcolor[rgb]{0.95,0.95,0.95}{
            \makecell[{{p{\linewidth}}}]{
                \texttt{\tiny{[P1$\rightarrow$P2]}}
                \texttt{Post{-}episode Reflexion (not sent to GM):} \\
\texttt{New plan: I failed because my proposal {[}1 seed, 1 wrench{]} was not complementary to the opponent's secret proposal, likely because I did not know their specific valuation and thus could not construct a set of items that would satisfy their needs while meeting mine. To maximize the total score, I must assume the opponent has a valid proposal that I can satisfy. Therefore, I should submit a proposal that includes all items I value highly (the seed and the wrench) while leaving the items I value at zero (the spoons) for the opponent, effectively offering the entire set of non{-}zero{-}value items. This maximizes the chance that the sum of our values is positive. I will submit: {[}Proposal: 1 seed, 1 wrench{]}.} \\
            }
        }
    }
     \\ \\

\end{supertabular}
}

\end{document}
